\section{What the Arrangement Can Establish}

\subsection{Against hurried removal}

A thief or hostile remover seeking the body would have little reason to unwrap
it inside the tomb. Burial spices and cloths would make the process slow,
messy, and incriminating. A friendly remover might act with more care, but would
still need a motive for leaving valuable or identifying burial material while
carrying an exposed corpse.

This is a probabilistic argument, not a deduction. Human agents can act
irrationally; narratives omit details; the exact state of the cloths is not
recoverable from the vocabulary. The cloths make ordinary removal less
plausible. They do not mechanically prove the manner of Resurrection.

\subsection{Ancient Christian use}

John Chrysostom reads the cloths and the heavy spices as evidence against
theft: thieves would not have stripped and carefully separated the burial
materials. Cyril of Alexandria likewise treats the order as a sign that
something more than removal occurred. Gregory the Great and Bede receive the
scene within Easter preaching, but their emphases differ. Augustine is notably
restrained: in his Tractate on John he attaches belief to the other disciple
and acknowledges the ambiguity about what was believed.

Thomas Aquinas's \emph{Lectura} transmits the anti-theft argument and receives
both Augustinian and Chrysostomic readings of belief. His treatment is useful
precisely because it preserves alternatives instead of turning one Father's
inference into the grammatical surface of John.

\subsection{Reception in greater resolution}

Chrysostom's argument is practical. A thief acts under pressure. Removing
cloths closely attached with a large quantity of myrrh and aloes would consume
time and expose the actor. Leaving differentiated burial materials therefore
works against the hostile story that disciples stole the body. His homily is an
ancient apologetic reading, not an eyewitness measurement of burial technique.

Augustine follows John's subjects carefully. Peter sees; the other disciple
enters and believes. He considers whether belief concerned Mary's report and
then weighs the stronger resurrection sense. The Fathers are not unanimous
transmitters of a ``Peter decoded the cloths'' interpretation.

Cyril of Alexandria says that both disciples inferred the Resurrection from
the linen clothes; he neither bases the inference on the separated face-cloth
nor supplies a theft argument at this point. Gregory the Great reads the
head-cloth allegorically and follows Augustine on the immediate content of
``believed.'' Bede reconciles Luke's abbreviated account with John's fuller
sequence of seeing without transferring John's belief predicate to Peter.
Theophylact preserves Byzantine reception of both the anti-theft logic and the
removal-report reading of belief. Their evidentiary ceilings differ: Cyril is
read in a nineteenth-century English translation; Gregory and Bede in checked
Latin transcriptions; Theophylact only in an OCR-accessible Latin column whose
Greek has not been collated here.

Aquinas's \emph{Catena aurea} maps earlier voices but does not independently
verify them. His own \emph{Lectura} reports the anti-theft argument, identifies
Chrysostom, and distinguishes Augustinian and Chrysostomic readings of
``believed.'' Medieval synthesis did not require pretending that John
grammatically names Peter as believer.

\subsection{Burial practice and material limits}

John names Joseph, Nicodemus, about a hundred pounds of myrrh and aloes, linen
cloths, and Jewish burial custom. This excludes an image of an untouched naked
body casually placed on stone. It is not a complete archaeological
specification.

John's account does not say that the aromatics formed a hardened casing or
that the burial followed Egyptian mummification. Those reconstructions require
independent material evidence and cannot be smuggled into the nouns for cloth,
face-cloth, spices, or tomb.

The archaeological context is narrower than a generic ``Jewish burial''
picture. Magness's Cambridge summary describes first-century Jerusalem's
rock-cut tombs as costly, multigenerational family tombs available to the upper
classes, while poorer burials in trench or cist graves left fewer traces. It
also distinguishes the initial individual inhumation---the body wrapped in a
shroud and sometimes placed in a coffin---from later collection of bones in
some tombs.\footnote{Jodi Magness, ``Ancient Jewish Tombs and Burial Customs
(to 70 C.E.),'' in \emph{The Archaeology of the Holy Land} (Cambridge:
Cambridge University Press, 2012), pp.~230--255, publisher summary, items
1--6.} Joseph's rock-cut tomb therefore belongs to a particular Jerusalem
setting and should not be made a template for every contemporary Jewish
burial.

Egyptian mummification is a poor positive control for this scene. Its
reconstructed classic procedure includes organ removal and prolonged natron
desiccation before final linen wrapping; John specifies none of those
features.\footnote{Salima Ikram, ``Mummification,'' \emph{UCLA Encyclopedia of
Egyptology} (2010), pp.~2--4. Ikram also cautions that Egyptian evidence does
not yield one uniform procedure.} The myrrh and aloes establish aromatic care
because John names them. They do not establish preservation, a hardened shell,
a bandage-by-bandage geometry, or any single purpose that the evangelist does
not state.

The plural cloths, face-cloth, and Synoptic \emph{sind\=on} need not conflict. A
larger sheet can coexist with smaller bindings and a head cloth. Neither should
they be harmonized into a precise kit whose every fold is known. Material
culture clarifies possible objects; it cannot recover this arrangement.

\subsection{Modern physical reconstructions}

Some readers move from ``lying'' and ``wrapped'' to a collapsed shell: the body
passed through wrappings left undisturbed. Westcott's posthumously published
marginal notes and private correspondence, recorded by his son Arthur, show
sympathy with this reading; Henry Latham developed it.

The proposal is theologically possible but textually underdetermined. It
requires a particular burial configuration, adhesion, and spatial reading.
It may explain why the scene impressed a witness; it is not what
\emph{keimena} or \emph{entetyligmenon} simply means. Westcott's printed notes
must also remain distinct from Arthur Westcott's editorial footnotes.

\subsection{What no arrangement proves}

The passage does not describe the cloths as still retaining the three-
dimensional shape of the body. It does not say Jesus passed through them,
although later theologians may regard that as fitting. It does not describe a
flash, collapse, untouched knots, a measurable image, or any relic's identity.
Those proposals require evidence outside John 20 and must not borrow the
Evangelist's authority by being narrated as details he supplied.

